\section*{Aufgabenstellung}
Anlaysieren Sie Registry Einträge bezüglich USB Devices , LAN und oder WLAN Informationen.

\section*{Lösung}

Zuerst wurde die SYTEM Registry aus dem Image extrahiert:
\begin{lstlisting}[language=bash]
root@Alpha-17:/home/luis/Programming/GitHub/University/DigitalForensic# fls -pr /dev/loop35p3 | grep 'SYSTEM'
r/r 37802-128-4:        Windows/Prefetch/SYSTEMPROPERTIESPROTECTION.EX-81A2FDE2.pf
r/r 109603-128-4:       Windows/Prefetch/SYSTEMSETTINGS.EXE-084BB8F9.pf
r/r 105387-128-4:       Windows/Prefetch/SYSTEMSETTINGSBROKER.EXE-8BBE2894.pf
r/r 101804-128-4:       Windows/System32/config/SYSTEM
r/r 101802-128-5:       Windows/System32/config/SYSTEM.LOG1
r/r 102064-128-5:       Windows/System32/config/SYSTEM.LOG2
root@Alpha-17:/home/luis/Programming/GitHub/University/DigitalForensic# icat /dev/loop35p3 101804-128-4 > /tmp/SYSTEM
root@Alpha-17:/home/luis/Programming/GitHub/University/DigitalForensic# file /tmp/SYSTEM
/tmp/SYSTEM: MS Windows registry file, NT/2000 or above
\end{lstlisting}

\subsection*{USB Einträge}

Im anschluss wurde nach USB Einträgen gesucht:

\begin{lstlisting}[language=bash]
regripper -p usbdevices -r /tmp/SYSTEM
regripper -p usbstor -r /tmp/SYSTEM
regripper -p devclass -r /tmp/SYSTEM
regripper -p mountdev -r /tmp/SYSTEM
\end{lstlisting}

Es wurden folgende 4 USB Geräte gefunden:

\subsubsection*{Gerät 1 – SanDisk 3.2Gen1}

\begin{tabular}{ll}
	\textbf{Hersteller / Modell}   & SanDisk 3.2Gen1 \\
	\textbf{VID / PID}             & \texttt{VID\_0781} / \texttt{PID\_5581} \\
	\textbf{Seriennummer}          & \texttt{010116ed7ac862...} (gekürzt) \\
	\textbf{Laufwerksbuchstabe}    & \texttt{D:} \\
	\textbf{Erste Installation}    & 2026-03-23 08:15:11Z \\
	\textbf{Letztes Entfernen}     & 2026-03-23 08:15:54Z \\
\end{tabular}

\medskip
\noindent
Das Gerät war lediglich ca.\ 43~Sekunden angeschlossen (zwischen \emph{Last Arrival} und \emph{Last Removal}), was auf ein schnelles Kopieren oder ein gezieltes kurzes Anschließen (testen, etc) hindeuten kann.

\subsubsection*{Gerät 2 – CnMemory USB Device}

\begin{tabular}{ll}
	\textbf{Hersteller / Modell}   & CnMemory \\
	\textbf{VID / PID}             & \texttt{VID\_090C} / \texttt{PID\_1000} \\
	\textbf{Seriennummer}          & \texttt{00040127001068} \\
	\textbf{Laufwerksbuchstabe}    & keiner (nur Volume-GUID vorhanden) \\
	\textbf{Erste Installation}    & 2026-03-23 07:52:45Z \\
	\textbf{Letztes Entfernen}     & 2026-03-23 08:15:06Z \\
\end{tabular}

\medskip
\noindent
Dem Gerät wurde kein Laufwerksbuchstabe zugewiesen; es ist lediglich eine
Volume-GUID eingetragen (\texttt{\{46308943-2685-11f1-a521-fa789cc76a14\}}).

\subsubsection*{Gerät 3 – SMI USB DISK (Instanz A)}

\begin{tabular}{ll}
	\textbf{Hersteller / Modell}   & SMI USB DISK \\
	\textbf{VID / PID}             & \texttt{VID\_090C} / \texttt{PID\_2000} \\
	\textbf{Seriennummer}          & \texttt{03161671000618} \\
	\textbf{Laufwerksbuchstabe}    & keiner (nur Volume-GUID vorhanden) \\
	\textbf{Erste Installation}    & 2026-02-23 13:03:25Z \\
	\textbf{Letztes Entfernen}     & nicht erfasst \\
\end{tabular}

\subsubsection*{Gerät 4 – SMI USB DISK (Instanz B)}

\begin{tabular}{ll}
	\textbf{Hersteller / Modell}   & SMI USB DISK \\
	\textbf{VID / PID}             & \texttt{VID\_090C} / \texttt{PID\_1000} \\
	\textbf{Seriennummer}          & \texttt{AA00000000013226} \\
	\textbf{Laufwerksbuchstabe}    & keiner \\
	\textbf{Erste Installation}    & 2026-02-16 13:42:54Z \\
	\textbf{Letztes Entfernen}     & nicht erfasst \\
\end{tabular}

\subsection*{LAN / WLAN Einträge}

Nun habe ich nach LAN und WLAN Einträgen gesucht:

\begin{lstlisting}[language=bash]
regripper -p nic2 -r /tmp/SYSTEM
regripper -p networklist -r /tmp/SOFTWARE
regripper -p ssid -r /tmp/SOFTWARE
\end{lstlisting}

ich habe hier die folgenden 3 Adapter gefunden:

\subsection*{Adapter 1 – \texttt{\{20b88bfa-bfd2-41ba-9779-1ea2a01031d6\}}}

\begin{tabular}{ll}
	\textbf{GUID}            & \texttt{\{20b88bfa-bfd2-41ba-9779-1ea2a01031d6\}} \\
	\textbf{LastWrite}       & 2026-02-16 13:42:56Z \\
	\textbf{DHCP aktiviert}  & Ja \\
	\textbf{IP-Adresse}      & nicht zugewiesen (kein DHCP-Lease) \\
	\textbf{DNS-Server}      & nicht konfiguriert \\
	\textbf{Domain}          & nicht konfiguriert \\
\end{tabular}

\medskip
\noindent
Der Adapter ist DHCP-fähig, hat jedoch keine IP-Adresse erhalten. Es liegt kein
aktiver Lease vor, was darauf hindeutet, dass zum Zeitpunkt des letzten Schreibens
keine Netzwerkverbindung bestand.

\subsection*{Adapter 2 – \texttt{\{64ac53af-0b3d-11f1-a51a-806e6f6e6963\}}}

\begin{tabular}{ll}
	\textbf{GUID}            & \texttt{\{64ac53af-0b3d-11f1-a51a-806e6f6e6963\}} \\
	\textbf{LastWrite}       & 2026-02-16 13:43:02Z \\
	\textbf{Konfiguration}   & keine Subkeys vorhanden \\
\end{tabular}

\medskip
\noindent
Für diesen Adapter sind keinerlei TCP/IP-Parameter hinterlegt. Es handelt sich
vermutlich um einen virtuellen oder nicht genutzten Adapter (z.\,B.\ ein
Loopback- oder VPN-Adapter).

\subsection*{Adapter 3 – \texttt{\{655f726f-6fea-44e7-a87e-2bea2176bcdb\}}}

\begin{tabular}{ll}
	\textbf{GUID}              & \texttt{\{655f726f-6fea-44e7-a87e-2bea2176bcdb\}} \\
	\textbf{LastWrite}         & 2026-03-23 06:53:35Z \\
	\textbf{DHCP aktiviert}    & Ja \\
	\textbf{IP-Adresse}        & \texttt{192.168.2.91} \\
	\textbf{Subnetzmaske}      & \texttt{255.255.255.0} (\texttt{/24}) \\
	\textbf{DHCP-Server}       & \texttt{192.168.2.1} \\
	\textbf{Standard-Gateway}  & \texttt{192.168.2.1} \\
	\textbf{DNS-Server}        & \texttt{192.168.2.1} \\
	\textbf{Domain}            & \texttt{fritz.box} \\
	\textbf{Lease erhalten}    & 2026-03-23 06:53:35Z \\
	\textbf{T1 (Renewal)}      & 2026-03-28 06:53:35Z \\
	\textbf{T2 (Rebinding)}    & 2026-04-01 00:53:35Z \\
	\textbf{Lease läuft ab}    & 2026-04-02 06:53:35Z \\
	\textbf{Lease-Dauer}       & 864.000~Sekunden (10~Tage) \\
\end{tabular}

\medskip
\noindent
Dieser Adapter war zum Zeitpunkt der Untersuchung aktiv mit einem Heimnetzwerk
verbunden. Die Domain \texttt{fritz.box} identifiziert den Router eindeutig als
eine \textbf{AVM FRITZ!Box}. Gateway und DNS-Server sind identisch
(\texttt{192.168.2.1}), was dem typischen Heimrouter-Setup entspricht.
Die Lease-Dauer von 10~Tagen sowie der Renewal-Zeitpunkt (T1) am 2026-03-28
deuten auf eine stabile, längerfristige Verbindung hin.
